\subsubsection{Using Samples for Scheduling}

It is known from equation \eqref{eq:time_duration} that the number of samples 
define duration. Assuming that the sampling rate remains fixed, it provides a useful 
metric for timestamping and scheduling MIDI events, whenever a melody is written. 
For example, given the sample rate ($f_s$), and the tempo in BPM, the time duration 
of one beat in samples is estimated using the following equation:

\begin{equation}
    t = \left\lfloor f_s \times  \frac{60}{\text{BPM}} \right\rfloor
\end{equation}

For example, if a music is played at 80 BPM, and the sample rate is \SI{44.1}{\kilo \hertz},

$$
    t = \left\lfloor 44100 \times \frac{60}{80} \right\rfloor = 33075
$$

This means that 33,075 samples are required to make one beat. However, the length of one 
music note is not limited to one beat.

Using the same tempo and sampling rate, when a note is scheduled at the 3rd beat, it needs to wait until
the first $33075 \times 2 = 66150$ samples are processed.

This makes it very useful during mixing all the tracks using samples, just to keep the timings
consistent without any interruption. It was initially thought of using \texttt{sleep\_for()} from
the C++ \texttt{std::chrono} library, which pauses an execution for a certain amount of time.
However, this introduces complexities which might require creating enormous amount of threads. Sometimes, they could halt
any high-priority callbacks which may cause bugs. But most importantly, sometimes \texttt{sleep\_for()}
can block for longer than the specified duration due to OS-related scheduling and resource contention 
\cite{cppreference_sleep_for}, which is not real-time friendly by any means. 

To further support this statement, audio hardware already provides the most accurate clock
in the system (\Cref{prompt:chronovssample}) \cite{openai_chatgpt}. Therefore, the audio callback 
timing using the number of samples provides an accurate scheduling of each MIDI events in the 
audio engine.